\begin{frame}{Literature Review (Paper 5)}
\begin{columns}[T]
  \begin{column}{0.25\textwidth}
    \begin{block}{Title of the paper \cite{ama26}}
      \tiny \textbf{AMA: Adaptive Memory via Multi-Agent Collaboration}\\[2pt] arXiv:2601.20352 (2026).\\[2pt] Problem: multi-agent systems need memory that can adapt to task context and uncertainty instead of relying on a fixed memory schema.
    \end{block}
    \vspace{0.12cm}
    \begin{block}{Main contributions}
      \tiny
      \begin{itemize}\setlength\itemsep{1pt}
      \item Introduces explicit memory lifecycle stages across formation and update.
\item Uses varying retrieval granularity and query routing.
\item Incorporates judgement, retry and conflict detection.
\item Supports refresh and targeted memory update when evidence is insufficient.
      \end{itemize}
    \end{block}
    \vspace{0.12cm}
    \begin{block}{Inputs and Outputs}
      \tiny Agent experience, task queries and evidence are transformed into adaptive memory updates and context-specific retrieval decisions.
    \end{block}
  \end{column}

  \begin{column}{0.4\textwidth}
    \begin{block}{Methodology}
    \tiny
    \begin{center}
    \textbf{Adaptive memory lifecycle}
    \end{center}
    \begin{itemize}\setlength\itemsep{1pt}
    \item Extract reusable experience from agent actions.
\item Route retrieval based on task query and memory granularity.
\item Judge reliability and sufficiency of evidence.
\item Refresh or revise memory when conflicts or gaps emerge.
    \end{itemize}
    \end{block}

    \begin{block}{Tools and Technology used}
      \tiny Multi-agent collaboration patterns; adaptive retrieval; lifecycle-based memory control; evidence-aware update logic.
    \end{block}
  \end{column}

  \begin{column}{0.25\textwidth}
    \begin{block}{Summary of Results}
      \tiny The paper highlights adaptive memory as a dynamic lifecycle, where agents decide when to retrieve, when to retry and when to revise memory based on evidence quality.
    \end{block}
    \vspace{0.12cm}
    \begin{block}{Major Findings}
      \tiny
      \begin{itemize}\setlength\itemsep{1pt}
      \item Memory must adapt to task context and uncertainty.
\item Retrieval should not be fixed at one granularity.
\item Judgement and refresh are necessary for durable memory quality.
      \end{itemize}
    \end{block}
    \begin{alertblock}{SYNAPSE Adapts}
      \tiny SYNAPSE adopts the lifecycle idea as formation $\rightarrow$ selective admission $\rightarrow$ judgement $\rightarrow$ conflict resolution $\rightarrow$ consolidation $\rightarrow$ evolution.
    \end{alertblock}
  \end{column}
  \end{columns}
\end{frame}
